%------------------------------------------------------
% Harvard-style one-page resume -- NGUYEN KHAC BAO
% Target: AI Engineer Intern -- Algorithm & Agent
% Compile: pdflatex CV_Nguyen_Khac_Bao_FPT_Algorithm_Agent_Intern.tex
%------------------------------------------------------

\documentclass[10pt,a4paper]{article}

\usepackage[T1]{fontenc}
\usepackage[utf8]{inputenc}
\usepackage[a4paper,top=0.38in,bottom=0.38in,left=0.52in,right=0.52in]{geometry}
\usepackage{enumitem}
\usepackage{hyperref}
\usepackage{titlesec}
\usepackage{microtype}
\usepackage{xcolor}

\hypersetup{
    colorlinks=true,
    urlcolor=blue,
    linkcolor=black,
    pdfborder={0 0 0}
}

\titleformat{\section}
    {\normalfont\large\bfseries}
    {}
    {0em}
    {}
    [\vspace{1pt}\titlerule]

\titlespacing*{\section}{0pt}{4pt}{2pt}

\setlength{\parindent}{0pt}
\setlength{\parskip}{0pt}
\pagestyle{empty}
\sloppy

\newenvironment{cvitems}{%
    \begin{itemize}[leftmargin=1.2em,itemsep=0pt,topsep=1pt,parsep=0pt]
}{%
    \end{itemize}
}

\newcommand{\entryhead}[2]{%
    \noindent\textbf{#1}\hfill #2\\[-2pt]
}

\newcommand{\entrysub}[1]{%
    \noindent\textit{#1}\\[-1pt]
}

\begin{document}
\fontsize{9.05}{10.0}\selectfont

\begin{center}
    {\Large\textbf{NGUYEN KHAC BAO}}\\[-1pt]
    AI Engineer Intern -- Algorithm \& Agent\\[2pt]
    Ho Chi Minh City, Vietnam
    \,\textbar\,
    \href{mailto:nguyenkhacbaowork564@gmail.com}{nguyenkhacbaowork564@gmail.com}
    \,\textbar\,
    (+84) 039 586 4429
    \,\textbar\,
    \href{https://github.com/NguyenKhacBao564}{github.com/NguyenKhacBao564}
\end{center}

\section{CAREER OBJECTIVE}
AI Engineer Intern candidate focused on building practical AI systems combining LLM/RAG agents, backend APIs, computer vision pipelines, and deployable inference services. Hands-on with Python, FastAPI, Gemini API, YOLOv8n, OpenCV, Docker, Cloud Run, and real-time AI demos.

\section{EDUCATION}
\entryhead{Posts and Telecommunications Institute of Technology (PTIT)}{2022 -- 2027}
\entrysub{Bachelor of Information Technology}
Relevant Coursework: Machine Learning, Deep Learning, Computer Vision.

\section{TECHNICAL SKILLS}
\textbf{AI / LLM / Agent:} RAG pipelines, Gemini API, embeddings, vector search, prompt engineering, chatbot workflow\\
\textbf{Backend / API:} Python, FastAPI, REST API, WebSocket, Streamlit, JSON/API integration\\
\textbf{ML / Computer Vision:} PyTorch, YOLOv8n, OpenCV, OCR, object detection, tracking, dataset QA\\
\textbf{Cloud / DevOps:} Docker, Linux, Git/GitHub, Google Cloud Run, Cloud SQL, Cloud Build, Vast.ai\\
\textbf{Data:} SQL basics, CSV/JSON processing, metadata tracking, evaluation reports

\section{LANGUAGES}
Vietnamese: Native \textbar\ English: TOEIC 550/990; technical reading proficiency for documentation, model cards, API references, and engineering articles.

\section{PROJECTS}
\entryhead{Vietnamese Legal RAG Chatbot}{2026}
\entrysub{Personal Project \textbar\ Python, FastAPI, Streamlit, Gemini API, Qdrant Cloud, Cloud SQL, Docker, Google Cloud Run \textbar\ Source: \href{https://github.com/NguyenKhacBao564/legal_RaG_Chatbot_VietNamese-System}{GitHub}}
\begin{cvitems}
    \item Built a cloud-ready RAG chatbot with FastAPI backend and Streamlit frontend for retrieval-grounded Vietnamese legal question answering.
    \item Integrated Gemini API, embeddings, and Qdrant vector search to retrieve relevant context before generating answers.
    \item Stored chat history with Cloud SQL/MySQL and configured API/database credentials through environment-based deployment.
    \item Designed extension points for async document indexing, OpenAI-compatible LLM serving, and future AI agent workflows.
\end{cvitems}

\entryhead{Real-Time Face Mask Compliance Detection System}{2026}
\entrysub{Personal Project \textbar\ Python, FastAPI, WebSocket, YOLOv8n, OpenCV, Docker, Google Cloud Run \textbar\ Source: \href{https://github.com/NguyenKhacBao564/Real-Time-Face-Mask-Compliance-Detection-System}{GitHub}}
\begin{cvitems}
    \item Trained a YOLOv8n 3-class mask compliance detector: correct\_mask / incorrect\_mask / no\_mask, achieving 0.970 mAP@50 and 0.749 mAP@50:95 on validation.
    \item Built a reproducible dataset pipeline to convert VOC/XML annotations to YOLO format, audit class distribution, and track edge cases such as lighting, occlusion, blur, and reflection.
    \item Developed a FastAPI REST + WebSocket inference service with browser webcam UI for real-time detection, violation counting, and test-frame capture.
    \item Containerized and deployed the service to Google Cloud Run; smoke-tested public REST/WSS inference with warm latency around 151 ms REST and 127 ms WSS.
\end{cvitems}

\entryhead{Traffic Violation Detection \& License Plate Recognition}{2026}
\entrysub{Personal Project \textbar\ Python, OpenCV, Ultralytics YOLO, ByteTrack, HyperLPR3, JSON \textbar\ Source: \href{https://github.com/NguyenKhacBao564/Traffic-Violation-Detection-and-License-Plate-Recognition}{GitHub}}
\begin{cvitems}
    \item Built an end-to-end fixed-camera CV pipeline for vehicle detection, ByteTrack tracking, traffic-light ROI state, stop-line crossing logic, plate OCR, and JSON evidence generation.
    \item Fine-tuned a YOLO plate detector on 8,598 CCPD images, achieving 0.989 precision, 0.998 recall, and 0.994 mAP50 on validation.
    \item Designed a reproducible license plate dataset workflow with annotation parsing, YOLO bbox conversion, OCR crops, negative samples, train/val/test split, and QA reports.
    \item Analyzed real-world issues such as low-light images, heavy blur, small/tilted plates, weather conditions, hard negatives, and missing/invalid YOLO boxes.
\end{cvitems}

\entryhead{AutoFall Labeler Tool}{2026}
\entrysub{Additional Project \textbar\ Python, Streamlit, OpenCV, YOLO-Pose, Pandas \textbar\ Source: \href{https://github.com/NguyenKhacBao564/autofall-labeler-tool}{GitHub}}
\begin{cvitems}
    \item Built a human-in-the-loop dataset assistant for fall-detection data with video frame extraction, YOLO-Pose suggestions, Streamlit review UI, QA reports, and reviewed-only export.
    \item Designed CSV-based source-of-truth records for metadata, annotations, keypoints, and review events to support reproducible dataset versioning.
    \item Reviewed and cleaned a 656-frame CCTV benchmark, producing 332 reviewed exportable samples with 90.2\% coarse fall/non-fall agreement.
\end{cvitems}

\section{ACTIVITIES}
\entryhead{AI Research Group -- University Club}{2025 -- Present}
\entrysub{Member}
Discussed Machine Learning, Deep Learning, Computer Vision, RAG/LLM applications, and AI system development.

\end{document}
